
\documentclass{article}
\usepackage{colortbl}
\usepackage{makecell}
\usepackage{multirow}
\usepackage{supertabular}

\begin{document}

\newcounter{utterance}

\centering \large Interaction Transcript for game `mmlu\_pro', experiment `mmlu\_pro\_computer\_science', episode 22 with Qwen3.6{-}27B.
\vspace{24pt}

{ \footnotesize  \setcounter{utterance}{1}
\setlength{\tabcolsep}{0pt}
\begin{supertabular}{c@{$\;$}|p{.15\linewidth}@{}p{.15\linewidth}p{.15\linewidth}p{.15\linewidth}p{.15\linewidth}p{.15\linewidth}}
    \# & \multicolumn{2}{c}{Player} && \multicolumn{2}{c}{Game Master} \\
    \hline

    \theutterance \stepcounter{utterance}  
    & & & \multicolumn{4}{p{0.6\linewidth}}{
        \cellcolor[rgb]{0.9,0.9,0.9}{
            \makecell[{{p{\linewidth}}}]{
                \texttt{\tiny{[P1$\langle$GM]}}
                \texttt{The following are multiple choice questions (with answers) about computer science. Answer only with "The answer is (X)" where X is the correct letter choice.} \\
\\ 
\texttt{Question:} \\
\texttt{A certain pipelined RISC machine has 8 general{-}purpose registers R0, R1, . . . , R7 and supports the following operations:} \\
\texttt{ADD Rs1, Rs2, Rd (Add Rs1 to Rs2 and put the sum in Rd)} \\
\texttt{MUL Rs1, Rs2, Rd (Multiply Rs1 by Rs2 and put the product in Rd)} \\
\texttt{An operation normally takes one cycle; however, an operation takes two cycles if it produces a result required by the immediately following operation in an operation sequence.} \\
\texttt{Consider the expression AB + ABC + BC, where variables A, B, C are located in registers R0, R1, R2. If the contents of these three registers must not be modified, what is the minimum number of clock cycles required for an operation sequence that computes the value of AB + ABC + BC?} \\
\texttt{Options:} \\
\texttt{A. 10} \\
\texttt{B. 6} \\
\texttt{C. 9} \\
\texttt{D. 5} \\
\texttt{E. 12} \\
\texttt{F. 4} \\
\texttt{G. 7} \\
\texttt{H. 8} \\
\texttt{I. N/A} \\
\texttt{J. N/A} \\
\texttt{The answer is (B)} \\
\\ 
\texttt{Question:} \\
\texttt{A compiler generates code for the following assignment statement.} \\
\texttt{G := (A + B) * C {-} (D + E) * F} \\
\texttt{The target machine has a single accumulator and a single{-}address instruction set consisting of instructions load, store, add, subtract, and multiply. For the arithmetic operations, the left operand is taken from the accumulator and the result appears in the accumulator. The smallest possible number of instructions in the resulting code is} \\
\texttt{Options:} \\
\texttt{A. 7} \\
\texttt{B. 5} \\
\texttt{C. 12} \\
\texttt{D. 4} \\
\texttt{E. 11} \\
\texttt{F. 6} \\
\texttt{G. 8} \\
\texttt{H. 10} \\
\texttt{I. 9} \\
\texttt{J. 3} \\
\texttt{The answer is (I)} \\
\\ 
\texttt{Question:} \\
\texttt{SHA{-}1 has a message digest of} \\
\texttt{Options:} \\
\texttt{A. 628 bits} \\
\texttt{B. 780 bits} \\
\texttt{C. 160 bits} \\
\texttt{D. 256 bits} \\
\texttt{E. 2048 bits} \\
\texttt{F. 820 bits} \\
\texttt{G. 128 bits} \\
\texttt{H. 512 bits} \\
\texttt{I. 1024 bits} \\
\texttt{J. 64 bits} \\
\texttt{The answer is (C)} \\
\\ 
\texttt{Question:} \\
\texttt{\_\_\_\_\_\_\_\_\_\_\_\_\_ can modify data on your system – so that your system doesn’t run correctly or you can no longer access specific data, or it may even ask for ransom in order to give your access.} \\
\texttt{Options:} \\
\texttt{A. Email Trojans} \\
\texttt{B. Backdoor Trojans} \\
\texttt{C. Rootkit Trojans} \\
\texttt{D. Banking Trojans} \\
\texttt{E. Trojan{-}Downloader} \\
\texttt{F. Spyware Trojans} \\
\texttt{G. Network Trojans} \\
\texttt{H. Worm Trojans} \\
\texttt{I. Ransom Trojan} \\
\texttt{J. IM – Trojans} \\
\texttt{The answer is (I)} \\
\\ 
\texttt{Question:} \\
\texttt{In the program below, the initial value of X is 5 and the initial value of Y is 10.} \\
\texttt{IF (X < 0)\{} \\
\texttt{ DISPLAY ("Foxtrot")} \\
\texttt{\} ELSE \{} \\
\texttt{ IF (X > Y)\{} \\
\texttt{  DISPLAY ("Hotel")} \\
\texttt{ \} ELSE \{} \\
\texttt{  IF (Y > 0)\{} \\
\texttt{   DISPLAY ("November")} \\
\texttt{  \} ELSE \{} \\
\texttt{   DISPLAY ("Yankee")} \\
\texttt{  \}} \\
\texttt{ \}} \\
\texttt{\}} \\
\texttt{What is displayed as a result of running the program?} \\
\texttt{Options:} \\
\texttt{A. November} \\
\texttt{B. Foxtrot} \\
\texttt{C. Zulu} \\
\texttt{D. Alpha} \\
\texttt{E. Charlie} \\
\texttt{F. Bravo} \\
\texttt{G. Yankee} \\
\texttt{H. Echo} \\
\texttt{I. Hotel} \\
\texttt{J. Delta} \\
\texttt{The answer is (A)} \\
\\ 
\texttt{Question:} \\
\texttt{You have just received a Basic Assembler Language programyou ran. As you turn to the execution page to check theprintout, you find that every fourth line of the program reads, for example, PSW =... 00004F17B4C4. What is the PSW and what information does it provide? In addition, what informationis provided in the three lines following the PSW statement?} \\
\texttt{Options:} \\
\texttt{A. The three lines following the PSW statement are error codes.} \\
\texttt{B. PSW is a password for system protection.} \\
\texttt{C. PSW is a post{-}system warning that flags potential security breaches, with the subsequent lines listing the affected files.} \\
\texttt{D. PSW denotes the Previous Software Version, with the following lines outlining the changes made in the current version.} \\
\texttt{E. PSW refers to Processor Status Word, indicating the processor's current tasks. The lines after detail the task queue.} \\
\texttt{F. PSW stands for Program Submission Workflow, showing the process flow, while the following lines provide the timestamps of each step completion.} \\
\texttt{G. PSW means Program Step Width, representing the step size in memory allocation; the next lines show memory usage statistics.} \\
\texttt{H. PSW is a programming language used by the system.} \\
\texttt{I. PSW is the identifier for the program start window, and the following lines are user credentials.} \\
\texttt{J. PSW stands for Program Status Word. It contains the value of the location counter, system protection information, and program interrupt status. The three lines following the PSW statement display the contents of the system's registers at the time of interruption.} \\
            }
        }
    }
     \\ \\

    \theutterance \stepcounter{utterance}  
    & \multicolumn{4}{p{0.6\linewidth}}{
        \cellcolor[rgb]{0.561,0.737,0.561}{
            \makecell[{{p{\linewidth}}}]{
                \texttt{\tiny{[P1$\rangle$GM]}}
                \texttt{The answer is (J)} \\
            }
        }
    }
    & & \\ \\

    \theutterance \stepcounter{utterance}  
    & & & \multicolumn{2}{p{0.3\linewidth}}{
        \cellcolor[rgb]{0.9,0.9,0.9}{
            \makecell[{{p{\linewidth}}}]{
                \texttt{\tiny{[GM$|$GM]}}
                \texttt{J} \\
            }
        }
    }
    & & \\ \\

    \theutterance \stepcounter{utterance}  
    & & & \multicolumn{2}{p{0.3\linewidth}}{
        \cellcolor[rgb]{0.9,0.9,0.9}{
            \makecell[{{p{\linewidth}}}]{
                \texttt{\tiny{[GM$|$GM]}}
                \texttt{J} \\
            }
        }
    }
    & & \\ \\

    \theutterance \stepcounter{utterance}  
    & & & \multicolumn{2}{p{0.3\linewidth}}{
        \cellcolor[rgb]{0.9,0.9,0.9}{
            \makecell[{{p{\linewidth}}}]{
                \texttt{\tiny{[GM$|$GM]}}
                \texttt{game\_result = WIN} \\
            }
        }
    }
    & & \\ \\

\end{supertabular}
}

\end{document}
